\odiseachapter{Sonetillos}{Las féminas de la Odisea}

Lo de la niñez es literal. Quien imagine a un poeta como animal solemne, embelesado en la contemplación del Parnaso, es que ha leído demasiada literatura francesa. La poesía es pasión, sí, pero también, y mucho, juego y risa. Para muestra un botón. A cuenta de unas correcciones finales que Juanma hacía a este texto, me arranqué con unos ripios endecasílabos, en los que el poeta aconseja:

\begin{versopropio}
Solo una coma, sí, eliminaría,\\
me opongo a que preceda a «milerrante»,\\
porque por ella el dativo ambulante,\\
sin duda a las Musas invitaría.
\stanzabreak
Cuidemos la perífrasis, cabales,\\
y quita coma después de «extravío»,\\
y antes de zarpar del puerto el navío,\\
arregla mayúsculas iniciales.
\stanzabreak
Ya nada más objeto, compañero,\\
Vayacondios, por el ponto mi Ulises,\\
que es vagabundo, mañoso y artero,
\stanzabreak
que es de todos los hombres el primero\\
y el último también, hombre de ardides,\\
como todo hombre, falso y verdadero.
\end{versopropio}

A los que Juanma respondía con la autoridad del patrón al marinero:

\begin{versocita}
Harto vivo el poema te ha salido\\
aunque de «milerrante», buen amigo, en otra afino\\
que es acusativo, no dativo
\end{versocita}

Y eso no es todo. Mi querido amigo Quique Ruiz\footnote{Cuando llegué a Puerto de Sagunto, tras mi infeliz niñez en Alba, me encontré un pueblo industrial, feo y acogedor. En Alba rondaban los cíclopes, en el Puerto, habitaban los feacios. Entre ellos, un grupo de poetas que marcaron mi adolescencia y que incluía a Paco Salinas y a Quique. Los dos fueron maestros y los dos me ofrecieron el regalo de su amistad y su sabiduría.} me mandaba, mientras yo completaba este capítulo, unos deliciosos sonetos (él los llama, pundonoroso, sonetillos), dedicados a las féminas de la Odisea.

\emph{Calipso}

\begin{versocita}
¿Y si la inmortalidad\\
que Calipso le ofrecía\\
con amatoria maestría\\
se hubiera hecho realidad?
\stanzabreak
¡Para el canto del poeta\\
qué desastre hubiera sido\\
convertir en mantenido\\
al héroe de las mil tretas!
\stanzabreak
Los aspirantes al trono\\
vivitos y coleando.\\
Kavafis buscando el tono
\stanzabreak
de los aedos divinos.\\
Y nosotros navegando\\
sin Ítaca, sin destinos.
\end{versocita}

\emph{NAUSICA}

I

\begin{versocita}
Cierto es que con la ayuda\\
de la oportuna Atenea,\\
envuelta en sales y arena\\
halló una forma desnuda.
\stanzabreak
Con grande desenvoltura\\
NAUSICA tuvo ganada\\
la última gran parada\\
del héroe en su andadura.
\stanzabreak
Su discreción y prudencia,\\
la contención de su ardor,\\
hicieron mella en la audiencia
\stanzabreak
de la Corte y Odiseo,\\
por su verbo de color,\\
cumplir pudo su deseo.
\end{versocita}

II

\begin{versocita}
Aunque resulte dudoso\\
que después de veinte años\\
volver a oler los rebaños\\
de Ítaca suene hermoso,
\stanzabreak
un legendario cantor\\
lo dejó dicho y contado:\\
por el tiempo que ha pasado\\
no hay argumento mejor.
\stanzabreak
Eso pidió y eso obtuvo\\
del espléndido rey feacio\\
que sin saberlo mantuvo
\stanzabreak
viva la más grande historia,\\
el más sublime prefacio\\
del que se tenga memoria.
\end{versocita}

III

\begin{versocita}
Hubiera sido imposible\\
sin NAUSICA el desenlace.\\
De su personaje nace\\
una corriente sensible.
\stanzabreak
Y a pesar de Poseidón,\\
embarcado en nave feacia\\
el Señor de la Falacia\\
finiquitó la excursión.
\stanzabreak
«¿Es mi casa? ¿Estoy seguro?\\
¿Es mi rostro el que ese espejo\\
me devuelve ajado y duro?»
\stanzabreak
Nadie lo reconoció.\\
Sólo su fiel perro viejo\\
movió el rabo y lo lamió.
\end{versocita}

\emph{PENÉLOPE}

I

\begin{versocita}
¿Quién ha sido el cabroncete\\
que me ha asignado el papel\\
de Penélope la fiel\\
en este eterno sainete?
\stanzabreak
Por no torcer el relato\\
aquí estoy dándole al hilo\\
a ver si me despabilo,\\
más aburrida que un gato.
\stanzabreak
¡Qué paciencia hay que tener\\
para pasar a la historia\\
como ejemplo de mujer!
\stanzabreak
No sé si valdrá la pena\\
gozar la futura gloria\\
a cambio de esta condena.
\end{versocita}

II

\begin{versocita}
Puesto que condena es\\
estar pegada a la rueca\\
sin descomponer la mueca\\
como figura de gres.
\stanzabreak
¡Que veinte años son nada\\
aguantando el desatino\\
de esperar a un libertino\\
al mando de una manada
\stanzabreak
no se lo cree ni el vate\\
que por ansia de aventura\\
mandó a mi rey al combate!
\stanzabreak
¡Rozando estoy la locura!\\
¡Ya mi carne se debate\\
en un mar de calentura!
\end{versocita}

\emph{CIRCE}

I

\begin{versocita}
Quien a mi reino se allega\\
sabe bien que en fiera muta\\
y que al poco se disputa\\
mis favores de gallega.
\stanzabreak
Hermes el entrometido\\
malbarató mi jugada,\\
completar la marranada\\
a Odiseo no he podido.
\stanzabreak
Si animal no quiere ser\\
será el festín de mi cama\\
antes del amanecer.
\stanzabreak
Con hechizo o sin hechizo\\
soy como dátil en rama,\\
soy cegador bebedizo.
\end{versocita}

II

\begin{versocita}
¡Lo prometí y lo cumplí!\\
¡Un año de gozadera!\\
El mortal tiene madera\\
de amante con pedigrí.
\stanzabreak
¡Qué felicidad Zeus mío!\\
¡Cada noche una verbena,\\
cada trago un quitapena,\\
cada día un desvarío!
\stanzabreak
Pero se acabó la fiesta.\\
Aunque me siente fatal\\
tengo que agachar la cresta
\stanzabreak
o se quiebra el engranaje\\
y se queda sin final\\
la narración de este viaje.
\end{versocita}

III

\begin{versocita}
Parta pues presto Odiseo\\
por el Ponto y que sus aguas\\
sean como mis enaguas\\
amorosas le deseo.
\stanzabreak
Le he soplado mis consejos:\\
Cuándo debe decidir\\
por dónde ir o no ir\\
si quiere llegar muy lejos.
\stanzabreak
Yo con este gesto espero\\
que el poeta me redima\\
y me trate con esmero.
\stanzabreak
De pasadas travesuras,\\
sin forzar mucho la rima,\\
ya he pagado las facturas.
\end{versocita}

Afortunadamente, en este universo, Ulises no aceptó la vida eterna y con ello nos regaló a todos una patria, quizás la única patria en la que vale la pena creer, la Ítaca que todos llevamos en el corazón.
